% applications/part_intro-DE.tex

\chapter*{Anwendungen}
\phantomsection
\addcontentsline{toc}{chapter}{Anwendungen}

Dieser Teil legt dar, wie die genehmigte AIM-Architektur in realen
Arbeitsabläufen zu ingenieurwissenschaftlichen Konsequenzen führt.

Die vorangehenden Teile haben die genehmigten Ingenieurbegriffe für
Alignment-Semantik, Zustand, Laufzeitauswertung und Realitätskopplung
entwickelt. Der Anwendungsteil zeigt, was diese Architektur in konkreten
Ingenieurkontexten ermöglicht.

Die Anwendungen folgen daher der Abfolge
\[
\begin{aligned}
&\text{genehmigte Ingenieurbegriffe}
\rightarrow
\text{Anwendungskontext}
\\
&\rightarrow
\text{ingenieurwissenschaftlicher Anwendungsfall}
\\
&\rightarrow
\substack{\text{erforderliche Laufzeit-/}\\\text{Realisierungsdienste}}
\rightarrow
\substack{\text{abgeleitete Repräsentation}\\\text{oder Entscheidungsunterstützung}}
\rightarrow
\text{ingenieurwissenschaftliche Konsequenz}.
\end{aligned}
\]

Anwendungskapitel verwenden kanonische Begriffe und werden nicht zu neuen
begrifflichen Heimaten für Alignment-Identität, pose2/pose3, Realisierung
oder Repräsentationsebenen.

\section*{Bezug zur AIM-Roadmap}

Innerhalb der Roadmap entspricht dieser Teil dem Einsatz der etablierten
Architektur in praktischen Ingenieurarbeitsabläufen.

Das Symbolbild behandelt die leitende Anwendungsfrage: Wie können
nachgelagerte Systeme kanonische Semantik verwenden, ohne die begriffliche
Zuständigkeit zu verlagern?

\begin{figure}[ht]
\centering
% figures/applications/icon_applications_downstream-DE.tex
\begin{tikzpicture}[font=\small, node distance=9mm and 10mm]
  \node[draw, rounded corners, fill=blue!8, minimum width=31mm, minimum height=10mm] (canonical) {kanonische Bedeutung};
  \node[draw, rounded corners, fill=green!8, right=of canonical, minimum width=29mm, minimum height=10mm] (services) {Laufzeitdienste};
  \node[draw, rounded corners, fill=orange!10, right=of services, minimum width=31mm, minimum height=10mm] (downstream) {nachgelagerte Systeme};
  \draw[-{Stealth[length=2.4mm]}, thick] (canonical) -- (services);
  \draw[-{Stealth[length=2.4mm]}, thick] (services) -- (downstream);
  \node[below=5mm of services, text width=67mm, align=center, text=gray!70!black]
    {Repräsentationen verwenden Bedeutung; sie übernehmen nicht deren Zuständigkeit.};
\end{tikzpicture}

\caption{Anwendungssymbol: Die kanonische Bedeutung bleibt der Repräsentation und externen Systemen vorgelagert.}
\label{fig:icon-applications-downstream}
\end{figure}

\begin{principle}[Anwendungsprinzip]
\label{princ:application-principle}
In AIM zeigen Anwendungen die architektonische Befähigung und die
ingenieurwissenschaftliche Konsequenz, während die begriffliche
Zuständigkeit bei Grundlagen, Zustand, Laufzeit und Realität verbleibt.
\end{principle}

\begin{summary}
Anwendungen in AIM sind konsequenzorientierte Darstellungen etablierter
Ingenieursemantik.

Sie zeigen, wie kanonische Begriffe Entscheidungsunterstützung,
Interoperabilität und operative Ingenieurergebnisse ermöglichen, ohne den
begrifflichen Kernel neu zu definieren.
\end{summary}

Nachdem diese Anwendungen entwickelt sind, kann der Diskussionsteil ihre
Folgerungen im historischen und methodischen Kontext einordnen.
